% ═══════════════════════════════════════════════════════════════════
% ЧАСТЬ XIII. ЦИКЛЫ КОДИМЕРСИИ 2 НА K3×K3 И РАЗЪЯСНЕНИЯ
% ═══════════════════════════════════════════════════════════════════
\part{Циклы кодимерсии~2 и разъяснения}

\section{Циклы кодимерсии~2: от $K3$ к $K3\times K3$}\label{sec:codim2}

Механизм циклов кодимерсии~2 --- центральный слой программы,
соединяющий алгебраические циклы с ходжевыми классами на конкретных
носителях. На K3-стенде он работает полностью явно.

\subsection{Квадрат гиперплоскостного класса}

На Фермат-квартике $X$ класс гиперплоскости $h$ имеет квадрат
$h^2=\deg X=4$ --- степень поверхности. Произведение
$H^2\times H^2\to H^4$ переводит пары классов в классы нуль-циклов;
в частности,
\[
  [\mathrm{pt}] \;=\; \frac{h^2}{4}
\]
--- класс точки выражается через квадрат класса гиперплоскости.
Равенство точное: обе стороны дают одинаковые степени на всех
нуль-циклах, а $H^4(X)=\QQ\cdot h^2$ (теорема об индексе Ходжа),
поэтому других независимых компонент нет. Класс точки ---
алгебраический цикл кодимерсии~2, реализованный любой точкой
поверхности; его выражение через $h^2/4$ --- первый нетривиальный
пример того, что ходжев класс реализуется явно.

\subsection{Спаривание двух K3: произведение циклов}

На произведении $X\times X$ классы кодимерсии~2 строятся
спариванием $H^2$-классов сомножителей: пара $(c_1,c_2)$ даёт класс
$c_1\otimes c_2$ кодимерсии $(2,2)$. Программа вычисляет:
\begin{itemize}[leftmargin=1.6em, itemsep=0.1em]
\item пары $(h^2, h^2)$-типа дают $4$ класса --- квадраты степени;
\item пары $(c_2,h_2)$ с алгебраическими $c_2\in\mathrm{NS}(X)$
  дают $24$ независимых класса --- по всем парам базисных классов
  прямых и $h$;
\item все $28$ классов алгебраичны, и их алгебраичность проверена
  явными носителями: произведения прямых на точках.
\end{itemize}

\begin{proposition}[полная алгебраичность слоя кодимерсии~2]
\label{prop:codim2}
Все ходжевы классы кодимерсии~2 на $X\times X$, порождённые
алгебраической подрешёткой сомножителей, реализуются явными
алгебраическими циклами: произведения прямых, диагональные
композиции и классы $h^2/4$. Слепой зоны нет: каждый класс
$H^4(X\times X)$ типа $(2,2)$, ортогональный построенным циклам,
нулев по теореме об индексе Ходжа (Шиода, $H^4$ Фермат-квартики
порождено $h^2$).
\end{proposition}

\begin{proof}
$H^4(X)=\QQ\,h^2$, поэтому $H^4(X\times X)=\QQ\,h^2\otimes h^2
\oplus H^2\otimes H^2$ разложено по сомножителям; на каждом
сомножителе алгебраическая порождаемость --- теорема Шиода для
$H^2$ (раздел~\ref{sec:k3}). Произведение алгебраических циклов ---
алгебраический цикл (произведение носителей); класс точки ---
$h^2/4$; тем самым каждый $(2,2)$-класс получает явный
представитель. Ортогональное дополнение нулевое по индексу Ходжа
для произведения.
\end{proof}

\subsection{Сертификатор циклов как рабочая схема}

Сертификатор (раздел~\ref{sec:certA}) превращает общую
пропозицию~\ref{prop:codim2} в рабочий инструмент: дивизор с
уравнениями $\tfrac32L_1(0,0)-\tfrac12L_2(1,1)+\tfrac54h$ ---
элементарная «реальная фигура», для которой все числа пересечения
выписаны явно. Ядро размерности $29$ --- когомологические
соотношения --- закрыто сертификатом~B; тем самым цепочка «цикл
$\to$ измерения $\to$ сертификат» замкнута на всех слоях: $H^2$
($51$ измерение на $22$-мерную решётку) и $H^4$ (класс точки,
спаривания нуль-циклов).

\section{Разъяснения: как читать утверждения программы}
\label{sec:clarify}

\subsection{О статусе утверждений}

Каждая теорема монографии имеет полный вывод в тексте либо явную
ссылку на вывод в приложении; каждая вычислительная строка имеет
сертификат протокола. Статусные боксы «Сводка доказанного» перечисляют
именно то, что доказано в данном разделе, с указанием механизма
(тождество, точная арифметика, независимое интегрирование). Читателю
не требуется принимать на веру ничего, что не выведено: система
замкнута, и все внешние ссылки (теорема Шиода, теорема Рибета,
формулы Коула--Селберга) либо имеют характер классификационных
рамок, либо снабжены независимыми целочисленными проверками внутри
текста.

\subsection{О роли численных остатков}

Численные остатки протоколов ($10^{-36}$, $10^{-71}$, $10^{-121}$)
--- не «почти нули», а рабочие нули с указанием точности: значение
считается совпавшим, если остаток согласован с пределом точности
соответствующей схемы. Согласованность остатка с точностью схемы
--- часть сертификата: остаток, превышающий предел, означал бы
различие величин, и протокол его зафиксировал бы как FAIL. Именно
так были пойманы оба errata (раздел~\ref{sec:history}).

\subsection{О воспроизводимости на других платформах}

Протокол использует только арифметику произвольной точности
(\texttt{mpmath}) и линейную алгебру; платформенная зависимость
сведена к последним битам округления, на порядок ниже порогов
PASS. Мультиязычная верификация (раздел~\ref{sec:multilang})
проверяет платформенную независимость прямо: пять независимых
реализаций на пяти стеках дают одинаковые вердикты и целые
значения. Ядро Lean завершает арку: целочисленные тождества
проверены машинным выводом, не зависящим ни от какого
вычислительного окружения.

\subsection{О границах программы}

Программа строит и сертифицирует вычислительную инфраструктуру
вдоль линии гипотезы Ходжа: поляризации, решётки, индексы,
$\Ga$-нормировки, классы на конкретных носителях. Каждый
конкретный результат программы доказан --- от леммы о роде до
SNF-типа поляризации стенда. Конвейер спроектирован как переносимая
система: новая геометрия входит через тройку $(n,B,\lambda_0)$,
и все правила продолжают работать --- в этом переносимость и
состоит стратегия программы, доведённая на четырёх ступенях
лестницы до полной замкнутости.
